\pdfbookmark[0]{Observations of Magnetic Variation. 1834}{variation_1834_october}
\begin{center}\itshape
Observations of Magnetic Variation in Göttingen and Leipzig\\
on 1 and 2 October 1834.
\end{center}
\begin{center}
\hrule\vspace{2mm}
\textsc{Poggendorff}. Annalen der Physik und Chemie. 1834. Vol.\ 33.
\vspace{2mm}\hrule\vspace{5mm}
\end{center}
The observations of magnetic variation on the agreed days that I mentioned in my article on the magnetic observatory here have since been made here on two more occasions, on 6 and 7 August and on 23 and 24 September. Some quite strong and remarkable anomalies occurred during the first session. It is all the more regrettable that chance circumstances prevented corresponding observations elsewhere. The September observations, on the other hand, were also made in full
\SourceLeaf{532}in Leipzig and Berlin and almost in full in Braunschweig. Half of them were made in Copenhagen as well, where observations were mistakenly made on 24 and 25 September instead of 23 and 24 September. Publishing these observations in full would be of less interest, however, since the course of variation on those two days was very regular, although several very small anomalies in the first 24 hours showed an admirable agreement at all four sites. It remains remarkable, however, that according to a newspaper report, a very strong display of the northern lights was seen in Glasgow on the evening of 23 September. One can therefore say with certainty that any effect it had on the magnetic needle did not reach as far as northern Germany.

Professor \textsc{Weber}’s presence in Leipzig meanwhile prompted us to set a few extra hours for simultaneous observations in Göttingen and Leipzig. On 1 and 2 October the times chosen were $7\tfrac12$ to $8\tfrac12$ in the morning, $12\tfrac12$ to $1\tfrac12$ at midday, and 8 to 10 in the evening. Apart from a few minor omissions, both places followed this schedule. My son, \textsc{Wilhelm Gauss}, observed at the magnetic observatory here, and Professor \textsc{Weber}, Professor \textsc{Möbius} and Dr \textsc{Thieme} observed in Leipzig. — — — One will note with pleasure the close agreement in the large movements on the evening of 1 October and in almost all the small ones. This shows unmistakably that they arise from forces that act over great distances, although they remain very puzzling at present. In Leipzig the anomalies were generally a little smaller than in Göttingen. The seat of the forces at work will therefore have been nearer Göttingen. I will only add that during part of those hours I myself made observations with a second instrument installed in the astronomical observatory here. I intend to describe the instrument more fully soon. These observations showed an almost perfect parallelism with those of the magnetic observatory here, in larger and smaller movements alike. A similar result was obtained on 23 and 24 September, as well as in many other trials. On several occasions the clocks at the two sites have thus been compared to within a small fraction of a minute of time by means of the magnetic phenomena alone. The same can be done using
\SourceLeaf{533}the larger movements on 1 and 2 October between Göttingen and Leipzig, where at both places the clocks differed by only a few seconds from local mean time.

These experiences now make even the small fluctuations of the magnetic needle that change over very short intervals of the greatest interest. One must hope that they too will be followed closely by observations at many places far apart. It will therefore be essential to make all observations at quite short intervals. Up to now we have observed from 5-minute mark to 5-minute mark. Yet even this interval is almost too long, and in future we intend always to determine the position of the magnetic needle from 3-minute mark to 3-minute mark on the agreed days. I should mention here that we used the procedure explained by the editor of these Annals (Vol.\ XXXII, pp.\ 569 to 572) only at first, and have long since replaced it with a slightly modified one. To obtain the position of the magnetic needle at a particular instant, we observe it at six different instants, always separated by the duration of one swing, so that the desired instant falls in the middle. In place of the exact duration of a swing we choose the nearest whole number of seconds, or more precisely of clock beats, for example $20''$ in place of $20''4$ at the magnetic observatory. The observations on 2 October for $8^{\mathrm h}\,15'$ in the evening therefore ran as follows.
\begin{center}
\begin{tabular}{rrr|r}
$8^{\mathrm h}$ & $14'$ & $10''$ & $672.6$\\
 & & $30$ & $672.3$\\
 & & $50$ & $671.3$\\
 & $15$ & $10$ & $671.8$\\
 & & $30$ & $669.9$\\
 & & $50$ & $670.8$
\end{tabular}
\end{center}
\Needspace{13\baselineskip}
This gives five means, which actually correspond to the times given alongside them.
\begin{center}
\begin{tabular}{rrr|r}
$8^{\mathrm h}$ & $14'$ & $20''$ & $672.45$\\
 & & $40$ & $671.80$\\
 & $15$ & $0$ & $671.55$\\
 & & $20$ & $670.85$\\
 & & $40$ & $670.35$
\end{tabular}\\[2mm]
\begin{tabular}{lrr}
and from them the mean for & $8^{\mathrm h}\,15'$ & $671.40$
\end{tabular}
\end{center}

\SourceLeaf{534}I deliberately chose this example, in which the needle showed rapid changes. Which stand out this clearly even from 20-second mark to 20-second mark. We have plenty of cases where a similar result occurs even at intervals half as long. I must reserve for another occasion the explanation of certain variations on this method of observation, which we often use, and the justification of the way of taking the mean, which was chosen after careful consideration. But I must mention, since the editor of these Annals did not explicitly say so in the passage cited above, that an essential condition for all these methods of observation to be admissible is that the needle first be brought as nearly to rest as possible, so that its swings amount to only a few scale divisions. Here at the magnetic observatory, it is rarely necessary to settle the needle in this way, or at least to do so again in the course of an observation. But anyone observing in a less favourable place must be sure to do so in the interval, as often as necessary, by the known means.

Shorter intervals now appear necessary, making the observations much more laborious than before, when a record was required only from hour to hour. There have therefore been several calls for somewhat fewer and shorter observation sessions in future. — — —

Göttingen, 5 November 1834.

